\documentclass[12pt, twocolumn]{report}

\usepackage{caption}
\usepackage{float}
\usepackage{hyperref}
\usepackage{graphicx}
\usepackage[nottoc, numbib]{tocbibind}
\usepackage{outlines}

\author{Philipp Liermann 20-475873-20\\Supervisor: Prof. Dr. Raad Bin Tareaf}
\date{\today}

\begin{document}

\title{
\begin{center}
	\includegraphics[width=8cm]{./images/logo.png}
\end{center}
\vspace{2cm}
Hardening MFA web applications to counter evolving transparent phishinging attack
vectors
\vspace{2cm}
\large }

\maketitle

\newpage
\tableofcontents

\newpage
\section{Problem statement}
These days most websites have implemented two or multi-factor authentication systems
to prevent malicious third-party actors from using stolen login credentials to
commit identity theft. Despite this, a new technique referred to as transparent
proxy phishing still provides malicious actors the means to bypass two-factor authentication
systems during social engineering based phishing attacks. This creates a
critical security vulnerability, as this technique can be used to gain unauthorized
access to otherwise secure systems. To resolve this issue, it is imperative to
develop and implement effective measures to counteract transparent phishing and
other methods of circumventing multi-factor authentication systems.

\section{Research question}
How can we protect multi factor authentication secured web applications from transparent
proxy phishing attacks? Based on this question the following sub questions
were derived:
\begin{itemize}
	\item What exactly are classical phishing attacks in general and how do they
	      differ from the new attack vector?

	\item How big is the threat imposed by this new type of phishing attack?

	\item Why are many organizations still not aware of this issue?

	\item What can be done to help developers counteract this type of security threat?
\end{itemize}

\section{Objective}
As credential theft and cyberfraud in general are still a growing problem in
the digital age, it is important to develop and implement effective measures to
counteract this threat. The objective of this thesis is to find and document
different strategy's that can help organizations and developers to mitigate
this new attack vector.

\section{Theoretical foundations \& current state of research}
Cyberfraud is a form of internet-based fraud, usually involving the use of false
identities and/or stolen information to illegally obtain money, property, or services.
Cyberfraud is an increasingly pervasive problem that is becoming increasingly
difficult to combat, as fraudsters become more sophisticated in their methods.
In 2020, cyberfraud was estimated to cost the global economy over \$6 trillion\cite{6trillion},
with the financial sector suffering the most damages. The social, economic and
reputational costs of cyberfraud can be incredibly damaging, and can range from
the loss of money, to identity theft, to the disruption of businesses.
Cyberfraud has become so pervasive that it is essential for businesses and
individuals to take measures to protect themselves from it. This includes using
strong passwords, using two-factor authentication, and staying up to date with
the latest security protocols. \\ \\ A phishing attack is a type of
cyberattack in which an attacker attempts to gain confidential information, such
as passwords, credit card numbers, or other sensitive information, by sending
emails or other messages disguised as legitimate entities. These messages often
include malicious links to faked login prompts that will steal the victims credentials
upon entering them. These attacks are becoming increasingly sophisticated and
difficult to recognize, making it important for everyone to remain vigilant and
take steps to protect against them. \\ \\ A HTTP reverse proxy is a type of
proxy server that retrieves resources on behalf of a client from one or more
servers. This type of proxy is sometimes referred to as a “gateway” or “tunneling”
proxy because it acts as a gateway for the traffic to and from the server. A
reverse proxy will typically receive a request from a client, then forward that
request to an appropriate server on the same network. It then retrieves the
response from the server and sends it back to the client. This type of proxy server
is most often used in enterprise networks to protect against malicious traffic,
to balance load between multiple servers, and to cache static content. \\ \\
Transparent proxy phishing is a new technique used by attackers to intercept
and steal multi-factor authentication (MFA) tokens from unsuspecting users. Instead
of copying HTML code from the original page that the attacker is trying to impersonate,
this new attack uses a HTTP reverse proxy to just redirect the users traffic
to the original page. The attacking proxy operator can view and modify all traffic
that is going through it while the victim sees a one by one copy of the original
login page. By doing so login credentials and 2FA tokens can be extracted
easily. \\ \\ TLS is a cryptographic protocol that provides end-to-end security
for data sent between a client and a server. It is widely used to secure web
traffic, email, and other types of data. TLS is the successor to SSL, and is
often referred to as SSL/TLS. If a webserver is using TLS its URL starts with the
well known https:// prefix. \\ \\ TLS Fingerprinting is a technique used to identify
the TLS implementation of a client or server by analyzing the handshake process.
This can be used to identify which software is used by client or server. Industry
standards for fingerprinting algorithms have existed for a long time. These
include: JA3, JA3N and the whole JA4+ family. \\ \\

Many scientific papers have been published on this topic, but there is still a
lack of information on how to protect against this new attack vector. This is why
it is important to find and document different strategies that can help organizations
and developers to mitigate this new attack vector.

\section{Research design}
This thesis will be conducted in a quantitative research design. By analyzing
existing literature and scientific papers on the topic, but also by running own
experiments in which open source reverse proxy phishing toolkits will be used to
setup attack simulations with the goal to find flaws in their attack implementation.
With the gained knowledge from this experiments this paper will outline easy
to follow strategies to protect web services from this threat.

\newpage
\chapter{Literature Review}
\section{Overview of Multi-Factor Authentication}

Multi-Factor Authentication (MFA) is a security system that requires more than
one method of authentication from independent categories of credentials to verify
the user's identity for a login or other transaction. The most common
categories are \cite{mfa}:
\begin{itemize}
	\item Something the user knows (e.g. a password)

	\item Something the user has (e.g. a smartphone)

	\item Something the user is (e.g. a fingerprint)
\end{itemize}
MFA is used to protect the user from unauthorized access to their accounts,
and is widely used in the financial and healthcare industries, as well as in government
and military applications. MFA is also used in consumer applications, such as online
banking and e-commerce. The use of MFA is growing rapidly, as more and more organizations
recognize the need for stronger security measures to protect their users and their
data. \\ \\ MFA is a critical component of a strong security posture, and is an
essential tool for protecting against a wide range of cyber threats, including
phishing, credential theft, and identity theft. MFA is also an important tool for
protecting against insider threats, as it can help to prevent unauthorized
access to sensitive data and systems. MFA is also an important tool for protecting
against the growing threat of cyberfraud, as it can help to prevent
unauthorized access to financial accounts and other sensitive information. \\ \\

\section{Phishing Attacks: Evolution and Impact}
Trace the evolution of phishing attacks from simple scams to sophisticated
reverse proxy phishing, including case studies or statistics to underline their
significance.

Phishing attacks are a type of cyberattack in which an attacker attempts to gain
confidential information, such as passwords, credit card numbers, or other
sensitive information, by sending emails or other messages disguised as legitimate
entities. These messages often include malicious links to faked login prompts
that will steal the victims credentials upon entering them. These attacks are becoming
increasingly sophisticated and difficult to recognize. Phishing has evolved
over time, from simple scams to sophisticated attacks that are difficult to detect.
In the early days of the internet, phishing attacks were relatively simple and
easy to recognize. However, as technology has advanced, so have phishing
attacks. Today, phishing attacks are often highly sophisticated and difficult
to detect, making them a significant threat to individuals and organizations.
\\ \\

\section{Existing Countermeasures}
There are a number of strategies and technologies that can be used to thwart phishing
attacks, including those targeting MFA systems. These include:
\begin{itemize}
	\item Geo-Location based QR-Code Authentication Scheme to Defeat Active Real-Time
	      Phishing Attack \cite{link1}

	\item Analyzing 2FA Phishing Attacks and Their Prevention Techniques \cite{link2}

	\item Secure authentication scheme to thwart RT MITM, CR MITM and malicious
	      browser extension based phishing attacks \cite{link3}
\end{itemize}
These strategies and technologies have been developed to help protect against
phishing attacks, and have been shown to be effective in many cases. However, they
are not foolproof, and there are still many ways in which attackers can bypass
them. For example, attackers can use reverse proxy phishing toolkits to
intercept and steal MFA tokens, and bypass MFA systems

The paper "Catching Transparent Phish: Analyzing and Detecting MITM Phishing Toolkits" \cite{kondracki2021catching} provides a detailed
analysis of the most used reverse proxy phishing toolkit and its potential to
bypass MFA systems. The authors demonstrate how those toolkits can be used to intercept
and steal MFA tokens, and propose a detection method based on a statistical
model that evaluates a combination of TLS fingerprinting and response timing
analysis. The authors provide an AI based solution for finding transparent phishing
toolkits in the wild, but only provide limited advice on how to detect client
connections of those toolkits on the recieving server side.

\newpage

\chapter{Methodology}
\section{Research Design}
Outline the research approach, whether qualitative, quantitative, or mixed
methods, justifying the choice based on your research questions.

\section{Data Collection}
Describe how data will be gathered, including any experimental setups, simulations,
or surveys planned.

\section{Analysis Techniques}
Detail the methods for analyzing the collected data, such as statistical analysis,
content analysis, or software testing methodologies.

\chapter{Experimentation}
\section{Simulation of Reverse Proxy Phishing Attacks}
Describe the setup for simulating attacks, including the tools and
environments used.

\section{Evaluation of MFA Vulnerabilities}
Explain how the simulations will be used to evaluate the vulnerabilities in
current MFA implementations.

\section{Testing of Countermeasures}
Outline the process for developing and testing new or improved countermeasures
against these attacks.

\chapter{Results}
\section{Findings from Simulations}
Present the data collected from the simulations, providing analysis and interpretation.

\section{Effectiveness of Current Defenses}
Assess the effectiveness of existing defenses against reverse proxy phishing based
on your findings.

\section{Proposed Solutions}
Introduce any new solutions or improvements to existing solutions developed through
your research.

\chapter{Discussion}
\section{Implications of Findings}
Discuss the broader implications of your findings for cybersecurity practices and
MFA implementation.

\section{Limitations and Future Research}
Acknowledge any limitations of your study and propose areas for future research.

\newpage
\bibliographystyle{plain}
\bibliography{refs}
\end{document}
